% CP-VI-0038
% target_chapter: VI/28 — Normalization
% source_unit_id: IMP-DL-D27460A55B-P007
% source: Downloads/theory-of-algebraic-geometry-5.tex L5353-L5535
% solution_provenance: SOURCE_EXPLICIT_SOLUTION

\subsection*{Problem 8. Normalization}

\textbf{Problem.}
A scheme is normal if all its local rings are integrally closed domains.
Let \(X\) be an integral scheme. For each open affine subset
\(U=\operatorname{Spec} A\) of \(X\), let \(\overline{A}\) be the integral
closure of \(A\) in its quotient field, and let
\[
\widetilde{U}=\operatorname{Spec}\overline{A}.
\]
Show that one can glue the schemes \(\widetilde{U}\) to obtain a normal
integral scheme \(\widetilde{X}\), called the normalization of \(X\).
Show that there is a morphism
\[
\pi:\widetilde{X}\longrightarrow X
\]
with the following universal property: for every normal integral scheme
\(Z\), and for every dominant morphism \(f:Z\to X\), the morphism \(f\)
factors uniquely through \(\widetilde{X}\).

\textbf{Solution.}

Let \(X\) be an integral scheme. If \(U=\operatorname{Spec}A\subseteq X\)
is affine, then \(A\) is an integral domain. Its function field is
\[
K=\operatorname{Frac}(A),
\]
and \(\overline{A}\) denotes the integral closure of \(A\) in \(K\).

The basic compatibility needed for gluing is the following localization
fact.

\medskip

\textbf{Lemma.}
If \(S\subseteq A\) is a multiplicative subset, then the integral closure
of \(S^{-1}A\) in \(\operatorname{Frac}(A)\) is
\[
S^{-1}\overline{A}.
\]

\textbf{Proof of the lemma.}
Every element of \(S^{-1}\overline{A}\) is integral over \(S^{-1}A\),
because integrality is preserved under localization.

Conversely, suppose \(x\in \operatorname{Frac}(A)\) is integral over
\(S^{-1}A\). Then \(x\) satisfies a monic equation
\[
x^n+\frac{a_1}{s_1}x^{n-1}+\cdots+\frac{a_n}{s_n}=0,
\qquad a_i\in A,\ s_i\in S.
\]
After multiplying by a suitable element of \(S\), one obtains that
\(sx\) is integral over \(A\) for some \(s\in S\). Hence
\[
sx\in \overline{A},
\]
and therefore
\[
x=\frac{sx}{s}\in S^{-1}\overline{A}.
\]
Thus the integral closure of \(S^{-1}A\) is \(S^{-1}\overline{A}\).
\(\square\)

\medskip

Now let \(D(g)\subseteq U=\operatorname{Spec}A\) be a distinguished open
subset. Then
\[
D(g)=\operatorname{Spec}A_g.
\]
By the lemma, the normalization of \(D(g)\) is
\[
\operatorname{Spec}\overline{A_g}
=
\operatorname{Spec}(\overline{A}_g).
\]
But this is exactly the inverse image of \(D(g)\) inside
\(\operatorname{Spec}\overline{A}\). Hence normalization is compatible
with restriction to distinguished open subsets.

Since affine open subsets of a scheme have a basis of distinguished
open subsets, these canonical identifications allow us to glue the affine
schemes
\[
\widetilde{U}=\operatorname{Spec}\overline{A}
\]
along overlaps. The cocycle condition is automatic because all
identifications come from the same localization property above.

Thus the schemes \(\widetilde{U}\) glue to a scheme \(\widetilde{X}\).

The morphisms
\[
A\longrightarrow \overline{A}
\]
induce morphisms
\[
\operatorname{Spec}\overline{A}\longrightarrow \operatorname{Spec}A.
\]
These are compatible on overlaps, so they glue to a morphism
\[
\pi:\widetilde{X}\longrightarrow X.
\]

We now check that \(\widetilde{X}\) is normal and integral. Normality is
local on affine open subsets. Each affine piece of \(\widetilde{X}\) is
of the form
\[
\operatorname{Spec}\overline{A},
\]
where \(\overline{A}\) is integrally closed in its fraction field.
Therefore every affine coordinate ring of \(\widetilde{X}\) is integrally
closed. Hence \(\widetilde{X}\) is normal.

Also, since \(X\) is integral, the rings \(A\) are domains and their
integral closures \(\overline{A}\subseteq \operatorname{Frac}(A)\) are
again domains. The affine pieces glue along nonempty open subsets in a
compatible way, so \(\widetilde{X}\) is integral.

It remains to prove the universal property.

Let \(Z\) be a normal integral scheme, and let
\[
f:Z\longrightarrow X
\]
be a dominant morphism. Since \(X\) and \(Z\) are integral and \(f\) is
dominant, \(f\) induces an inclusion of function fields
\[
K(X)\hookrightarrow K(Z).
\]

Let \(U=\operatorname{Spec}A\subseteq X\) be affine, and let
\(W=\operatorname{Spec}B\subseteq f^{-1}(U)\) be affine. Since \(Z\) is
normal, \(B\) is an integrally closed domain.

The morphism \(f\) gives a ring homomorphism
\[
A\longrightarrow B.
\]
Now take any element \(\alpha\in \overline{A}\). By definition,
\(\alpha\) is integral over \(A\). Its image in \(K(Z)\) is therefore
integral over the image of \(A\), hence integral over \(B\). Since \(B\)
is integrally closed in its fraction field, the image of \(\alpha\)
actually lies in \(B\).

Therefore the map \(A\to B\) extends uniquely to a map
\[
\overline{A}\longrightarrow B.
\]
Equivalently, the morphism
\[
W\longrightarrow U
\]
lifts uniquely to a morphism
\[
W\longrightarrow \widetilde{U}.
\]

These local lifts are compatible on overlaps, because they are all
defined by the same induced map on function fields. Hence they glue to a
global morphism
\[
\widetilde{f}:Z\longrightarrow \widetilde{X}
\]
such that
\[
f=\pi\circ \widetilde{f}.
\]

Finally, uniqueness follows because on each affine open
\(\widetilde{U}=\operatorname{Spec}\overline{A}\), the map
\[
\overline{A}\longrightarrow \Gamma(W,\mathcal{O}_Z)
\]
is forced by the induced map of function fields
\[
K(X)\hookrightarrow K(Z).
\]
Thus there is at most one such lift.

Therefore \(\widetilde{X}\) is the normalization of \(X\), and
\(\pi:\widetilde{X}\to X\) satisfies the desired universal property.
